\chapter{Software organization, data schemas, and reproducibility records}
\label{ch:software}
\chapterpromise{How the implementation separates observation-model calculations from partition recursions, what the test runs established, and which implementation corrections remain incomplete.}

\section{Implemented organization}
The inspected package contains 44 modules and 9,143 source lines. Its principal components are
observation-family classes, partition dynamic programming, shared-changepoint and grouped-model
wrappers, nonconjugate numerical methods, posterior prediction, changepoint metrics, dataset
loaders, experiment scripts, and plotting functions. Observation-family modules populate arrays of
segment marginal likelihoods and posterior moments. The partition modules combine those arrays
with prior factors. Analysis modules return posterior summaries, a joint MAP partition, predictive
quantities, and result metadata.

\begin{figure}[H]
\centering
\resizebox{0.96\textwidth}{!}{
\begin{tikzpicture}[node distance=7mm and 8mm]
\node[primarybox,text width=23mm,double] (raw) {Raw data + hashes};
\node[secondarybox,right=of raw,text width=24mm] (load) {Loaders + schemas};
\node[primarybox,right=of load,text width=24mm] (model) {Block and DP modules};
\node[primarybox,right=of model,text width=23mm] (exp) {Experiment registry};
\node[successbox,right=of exp,text width=23mm,double] (side) {Result sidecars};
\draw[arrPrimary] (raw)--(load);\draw[arrPrimary] (load)--(model);\draw[arrPrimary] (model)--(exp);\draw[arrPrimary] (exp)--(side);
\node[secondarybox,below=10mm of side,text width=25mm] (assets) {Figures and tables};
\node[successbox,left=of assets,text width=25mm] (book) {Book and papers};
\draw[arrPrimary] (side)--(assets);\draw[arrPrimary] (assets)--(book);
\draw[arrSecondary,dashed,bend right=25] (book.west) to node[below,font=\scriptsize]{manifest verifies lineage}(raw.south);
\end{tikzpicture}
}
\caption{Relation among input data, model configuration, numerical calculation, returned results,
and generated figures or tables. A portable release should replace absolute author paths with data
identifiers and record data, configuration, environment, and output hashes.}
\label{fig:repository-provenance}
\end{figure}

\section{Observed test outcomes}
A clean Phase~0 run reported 174 passing tests and five skips caused by optional dependencies. The
historical Phase~1 bounded verification collected 179 tests: 157 passed, five were skipped because
of optional dependencies, and 17 nonconjugate or grouped-model cases did not complete. An
independent Phase~6 rerun collected the same 179 tests: 173 passed, five were skipped, one EP
logistic-normal test exceeded a 20-second per-test cap, and none failed. The remaining timeout is
neither a pass nor a failure, and no package code was modified for this verification.
\taskid{CODE-BB-012} specifies termination analysis and profiling for that remaining case. A source comparison found 110 scoped scientific files byte-identical to the fresh
extraction.

\begin{table}[H]
\centering
\begin{tabular}{@{}lrrrrr@{}}
\toprule
Verification context & Collected & Passed & Skipped & Unresolved & Failed\\ \midrule
Phase 0 package baseline & 179 & 174 & 5 & 0 & 0\\
Phase 1 bounded profile & 179 & 157 & 5 & 17 & 0\\
Phase 6 bounded rerun & 179 & 173 & 5 & 1 & 0\\
\bottomrule
\end{tabular}
\caption{Observed test outcomes. The rows use different execution limits and are reported
separately.}
\label{tab:test-status}
\end{table}

\section{Mathematical invariants checked in software}
\begin{enumerate}
\item Sum-product and max-sum recursions use the same observation-model and prior contributions but
different algebraic operations.
\item A boundary hazard is applied only at interior endpoints; for small $n$, the partition
normalizer is compared with exhaustive enumeration.
\item A latent-group optimization result should return the final value in its recorded objective
sequence, and restart ranking should use that value.
\item Comparisons among nonconjugate numerical methods must use the same admissible segments and the
same coordinate indexing.
\item Posterior prediction dispatches according to the fitted observation family and rejects
unsupported coordinates unless an extrapolation model is specified.
\item Boundary matching first maximizes the number of matches and then minimizes total distance;
matched-boundary MAE is \NA\ when no pair is matched.
\item Result metadata distinguish executed results, results valid for a stated analysis, results
excluded from a stated analysis, and planned calculations.
\end{enumerate}

\section{Reproduction commands and portable metadata}
The archived report includes editable-installation commands and synthetic or real-data entry
points. They identify intended calculations but do not establish that every external data source
is currently available. A reproducible run should record the package version, environment, random
seeds, model configuration, input hashes, cache history, and output metadata. Absolute paths in
archived result files are historical information and should be replaced by portable identifiers in
a release.

\begin{itemize}[leftmargin=1.7em,itemsep=2pt]
\item \path{python -m venv .venv}
\item \path{source .venv/bin/activate}
\item \path{pip install -e ".[all]"}
\item \path{pytest -q --durations=25}
\item \path{python -m bayesbreak.experiments.synthetic --all --out artifacts/synthetic}
\item \path{python -m bayesbreak.experiments.realdata --dataset <name> --out artifacts/realdata}
\end{itemize}

\section{Additional executed diagnostic outputs}
The archive also contains mixture-discovery, signal-to-noise sensitivity, multivariate
shared-changepoint, segment-count selection, and missing-data diagnostics. These outputs were
executed, but they are not required to establish the main results in this volume. Chapter~\ref{ch:synthetic}
uses the subset needed to describe boundary calibration, generality across observation models,
latent-group behavior, and finite-range timing.
